\ssr{ВВЕДЕНИЕ}
%<цель и задачи; можно вводное слово к работе>

Цель работы -- разработка и сравнительный анализ последовательного и параллельного алгоритмов.\\

Задачи, необходимые для достижения поставленной цели:
\begin{enumerate}
	\item описать последовательный алгоритм решения задачи по варианту и его параллельную версию;
	\item реализовать оба алгоритма;
	\item выполнить сравнительный анализ зависимостей времени решения задач от размерности входа для реализации последовательного алгоритма и для реализации модифицированного алгоритма, запущенной с единственным вспомогательным потоком;
	\item выполнить сравнительный анализ зависимостей времени решения задач от размерности входа для реализации модифицированного алгоритма при k вспомогательных потоках;
	\item сформулировать рекомендацию о выборе k для решения задачи на ЭВМ архитектуры.
\end{enumerate}

%Описать последовательный алгоритм решения задачи (по варианту).
%Разработать параллельную версию алгоритма.
%Реализовать обе версии алгоритма.
%Выполнить сравнительный анализ зависимостей времени решения задач от размерности входа для реализации последовательного алгоритма и для реализации модифицированного алгоритма, запущенной с единственным вспомогательным (рабочим) потоком.
%Выполнить сравнительный анализ зависимостей времени решения задач от размерности входа для реализации модифицированного алгоритма при k вспомогательных (рабочих) потоках, k принимает значения 1, 2, 4, ..., 8*q, где q -- количество логических ядер процессора Вашей ЭВМ.
%Сформулировать рекомендацию о выборе k для решения задачи на ЭВМ выбранной автором работы архитектуры.


%Дан граф в формате graphviz. Входные данные считываются из файла. Граф не полносвязный, количество вершин не менее 3000, из/в вершину ведут, вероятнее всего, от 0 до 15 дуг.
%
%Разрешено использовать только нативные потоки. Запрещены sleep, конструкции вида parallel for, стандартные пул потоков и менеджеры задач (в том числе планировщики и распределители решений): требуется распределять задачи между рабочими потоками, создав потоки вручную, из-под кода.

%

%В части 1 отчёта описать основные положения базового алгоритма расчёта, согласно варианту. Распараллелить базовый алгоритм --- разработать параллельный алгоритм расчёта, в части 1 описать его основные положения. Обосновать необходимость использования средств синхронизации --- мьютексов, семафоров --- или отсутствие оной.

%В части 2 описать схему последовательного алгоритма расчёта, схему алгоритма работы главного потока, который создаёт и запускает вспомогательные потоки, и по одной схеме алгоритма на каждый вид вспомогательного потока. Используемые структуры данных, в том числе мьютекс и/или семафор, должны быть описаны в данной части. Обоснуйте необходимость шага объединения промежуточных значений (или отсутствие таковой) и причины и место (в алгоритме) использования примитивов синхронизации.

%Реализовать описанный алгоритм последовательного расчёта и разработанный алгоритм параллельного расчёта. Использование стандартного пула потоков запрещено: требуется вручную создавать каждый поток, назначать ему задание и запускать его. В части 3 в числе используемых инструментов разработки указать используемые библиотеки и используемые реализации описанных в части 2 структур данных, в том числе мьютекса и/или семафора в случае, если они используются (всё со ссылками на официальную документацию). Также описать и проиллюстрировать необходимый порядок работы с сущностью вспомогательного потока, согласно используемой библиотеке. Для библиотек, стандартных структур данных и типа/класса потока необходимо дать ссылки на документацию к ним. Выполнить тестирование реализаций алгоритмов.

%В части 4 указать характеристики ЭВМ, на которой будут проведены замеры времени выполнения расчётов, включая число физических ядер центрального процессора и число логических ядер центрального процессора. Описать постановку эксперимента: варьирование размеров входов и параметров (при наличии параметров в алгоритме). Выполнить замеры времени выполнения расчёта для различных значений линейного(ых) размера(ов) входных данных, каждое приводимое в отчёте значение времени расчёта должно быть получено как усреднённое время серии расчётов --- серии идентичных экспериментов. Привести иллюстрацию сравнения времени расчёта реализацией последовательного алгоритма и времени расчёта реализацией параллельного алгоритма с одним вспомогательным потоком, интерпретировать полученные результаты. Также привести результаты сравнения времени выполнения реализации параллельного расчёта с варьирующимся числом вспомогательных потоков q и интерпретировать полученные результаты для k={1,2,4,8,16,…,8*q}, где q --- количество логических ядер на ЭВМ, на которой выполняются замеры времени. Дать рекомендацию о выборе значения k для данной ЭВМ и для предложенного алгоритма параллельного расчёта.

\clearpage
